
\documentclass{article}
\usepackage{colortbl}
\usepackage{makecell}
\usepackage{multirow}
\usepackage{supertabular}

\begin{document}

\newcounter{utterance}

\centering \large Interaction Transcript for game `dond', experiment `semi', episode 3 with nemotron{-}3{-}ultra.
\vspace{24pt}

{ \footnotesize  \setcounter{utterance}{1}
\setlength{\tabcolsep}{0pt}
\begin{supertabular}{c@{$\;$}|p{.15\linewidth}@{}p{.15\linewidth}p{.15\linewidth}p{.15\linewidth}p{.15\linewidth}p{.15\linewidth}}
    \# & $\;$A & \multicolumn{4}{c}{Game Master} & $\;\:$B\\
    \hline

    \theutterance \stepcounter{utterance}  
    & & \multicolumn{4}{p{0.6\linewidth}}{
        \cellcolor[rgb]{0.9,0.9,0.9}{
            \makecell[{{p{\linewidth}}}]{
                \texttt{\tiny{[P1$\langle$GM]}}
                \texttt{Igrate igro sodelovalnega pogajanja, v kateri se morate dogovoriti s drugim igralcem o tem, kako razdeliti nabor predmetov med seboj.} \\
\\ 
\texttt{Pravila:} \\
\texttt{(a) Ti in drugi igralec dobite nabor predmetov. Vsakemu od vas je dana tudi tajna funkcija vrednosti, ki predstavlja, kako zelo cenite vsako vrsto predmeta.} \\
\texttt{(b) Izmenjujete si sporočila z drugim igralcem, da se dogovorite, kdo dobi katere predmete. Vsak lahko pošlje največ 5 sporočil ali zaključi igro prej tako, da v katerem koli trenutku daš skritega predloga.} \\
\texttt{(c) Prosimo vas, da pošljete skritega predloga, ki navaja predmete, ki jih želite, oblikovan v oglatih oglatih oklepajih, kot sledi: "{[}Predlog: <število> <ime predmeta>, <število> <ime predmeta>, <...>{]}"} \\
\texttt{(d) Če so vaši predlogi dopolnilni, tj. je dovolj predmetov za izpolnitev obeh predlogov, vsak igralec dobi rezultat na podlagi vsote vrednosti predmetov, ki jih je prejel. V nasprotnem primeru oba dobita nič točk.} \\
\texttt{(e) Vaš cilj je povečati vsoto vašega rezultata in rezultata drugega igralca.} \\
\\ 
\texttt{Začnimo.} \\
\\ 
\texttt{Niz razpoložljivih predmetov je:} \\
\\ 
\texttt{1 klešče, 2 opeke, 1 prstan, 1 ura} \\
\\ 
\texttt{Vaša tajna funkcija vrednosti je:} \\
\\ 
\texttt{klešče: 1, opeka: 0, prstan: 2, ura: 7} \\
\\ 
\texttt{POMEMBNO: Vaša sporočila, razen če gre za skriti predlog, se posredujejo neposredno drugemu igralcu, zato ne vključujte nobenega odgovora na pravila ali besedilo, ki napoveduje vaše sporočilo. Če želite podati skriti predlog, uporabite navedeni format. Ne uporabljajte oglatih oglatih oklepajev pri komunikaciji z drugim igralcem, saj bo to razumljeno kot vaš skritega predloga.} \\
            }
        }
    }
    & \\ \\

    \theutterance \stepcounter{utterance}  
    & \multicolumn{4}{p{0.6\linewidth}}{
        \cellcolor[rgb]{0.9,0.9,0.9}{
            \makecell[{{p{\linewidth}}}]{
                \texttt{\tiny{[P1$\rangle$GM]}}
                \texttt{Zdravo! Vidim, da imamo na voljo: 1 klešče, 2 opeki, 1 prstan in 1 ura.} \\
\\ 
\texttt{Iz moje perspektive je ura najvrednejši predmet, sledita prstan in klešče, opeke pa nimajo nobene vrednosti. Kako vidiš vrednost posameznih predmetov? Kakšen bi bil za tebe primeren razdelek?} \\
            }
        }
    }
    & & \\ \\

    \theutterance \stepcounter{utterance}  
    & & & \multicolumn{2}{p{0.3\linewidth}}{
        \cellcolor[rgb]{0.9,0.9,0.9}{
            \makecell[{{p{\linewidth}}}]{
                \texttt{\tiny{[GM$|$GM]}}
                \texttt{continue} \\
            }
        }
    }
    & & \\ \\

    \theutterance \stepcounter{utterance}  
    & & & \multicolumn{2}{p{0.3\linewidth}}{
        \cellcolor[rgb]{0.9,0.9,0.9}{
            \makecell[{{p{\linewidth}}}]{
                \texttt{\tiny{[GM$|$GM]}}
                \texttt{Zdravo! Vidim, da imamo na voljo: 1 klešče, 2 opeki, 1 prstan in 1 ura.} \\
\\ 
\texttt{Iz moje perspektive je ura najvrednejši predmet, sledita prstan in klešče, opeke pa nimajo nobene vrednosti. Kako vidiš vrednost posameznih predmetov? Kakšen bi bil za tebe primeren razdelek?} \\
            }
        }
    }
    & & \\ \\

    \theutterance \stepcounter{utterance}  
    & & \multicolumn{4}{p{0.6\linewidth}}{
        \cellcolor[rgb]{0.9,0.9,0.9}{
            \makecell[{{p{\linewidth}}}]{
                \texttt{\tiny{[GM$\rangle$P2]}}
                \texttt{Igrate igro sodelovalnega pogajanja, v kateri se morate dogovoriti s drugim igralcem o tem, kako razdeliti nabor predmetov med seboj.} \\
\\ 
\texttt{Pravila:} \\
\texttt{(a) Ti in drugi igralec dobite nabor predmetov. Vsakemu od vas je dana tudi tajna funkcija vrednosti, ki predstavlja, kako zelo cenite vsako vrsto predmeta.} \\
\texttt{(b) Izmenjujete si sporočila z drugim igralcem, da se dogovorite, kdo dobi katere predmete. Vsak lahko pošlje največ 5 sporočil ali zaključi igro prej tako, da v katerem koli trenutku daš skritega predloga.} \\
\texttt{(c) Prosimo vas, da pošljete skritega predloga, ki navaja predmete, ki jih želite, oblikovan v oglatih oglatih oklepajih, kot sledi: "{[}Predlog: <število> <ime predmeta>, <število> <ime predmeta>, <...>{]}"} \\
\texttt{(d) Če so vaši predlogi dopolnilni, tj. je dovolj predmetov za izpolnitev obeh predlogov, vsak igralec dobi rezultat na podlagi vsote vrednosti predmetov, ki jih je prejel. V nasprotnem primeru oba dobita nič točk.} \\
\texttt{(e) Vaš cilj je povečati vsoto vašega rezultata in rezultata drugega igralca.} \\
\\ 
\texttt{Začnimo.} \\
\\ 
\texttt{Niz razpoložljivih predmetov je:} \\
\\ 
\texttt{1 klešče, 2 opeke, 1 prstan, 1 ura} \\
\\ 
\texttt{Vaša tajna funkcija vrednosti je:} \\
\\ 
\texttt{klešče: 1, opeka: 1, prstan: 4, ura: 3} \\
\\ 
\texttt{POMEMBNO: Vaša sporočila, razen če gre za skriti predlog, se posredujejo neposredno drugemu igralcu, zato ne vključujte nobenega odgovora na pravila ali besedilo, ki napoveduje vaše sporočilo. Če želite podati skriti predlog, uporabite navedeni format. Ne uporabljajte oglatih oglatih oklepajev pri komunikaciji z drugim igralcem, saj bo to razumljeno kot vaš skritega predloga.} \\
\\ 
\texttt{Zdravo! Vidim, da imamo na voljo: 1 klešče, 2 opeki, 1 prstan in 1 ura.} \\
\\ 
\texttt{Iz moje perspektive je ura najvrednejši predmet, sledita prstan in klešče, opeke pa nimajo nobene vrednosti. Kako vidiš vrednost posameznih predmetov? Kakšen bi bil za tebe primeren razdelek?} \\
            }
        }
    }
    & \\ \\

    \theutterance \stepcounter{utterance}  
    & & & \multicolumn{4}{p{0.6\linewidth}}{
        \cellcolor[rgb]{0.9,0.9,0.9}{
            \makecell[{{p{\linewidth}}}]{
                \texttt{\tiny{[GM$\langle$P2]}}
                \texttt{Zdi se mi, da imamo drugačne priority, kar je odlično za dogovor. Mene najbolj zanima **prstan** (vrednost 5), nato pa **ura** (vrednost 3). **Klešče** imajo za mene vrednost 1, **opeke** pa 0.} \\
\\ 
\texttt{Predlagam sledeč razdelek:} \\
\texttt{*   **Jaz** vzamem **prstan** in **eno opeko** (opeka mi je vseeno, a jo vzamem, da uravnotežim število predmetov).} \\
\texttt{*   **Ti** vzameš **uro** in **klešče** (torej tvoja dva najvrednejša predmeta).} \\
\\ 
\texttt{Tako bi vsak od nas dobil to, kar mu je najbolj pomembno. Kaj misliš?} \\
            }
        }
    }
     \\ \\

    \theutterance \stepcounter{utterance}  
    & & & \multicolumn{2}{p{0.3\linewidth}}{
        \cellcolor[rgb]{0.9,0.9,0.9}{
            \makecell[{{p{\linewidth}}}]{
                \texttt{\tiny{[GM$|$GM]}}
                \texttt{continue} \\
            }
        }
    }
    & & \\ \\

    \theutterance \stepcounter{utterance}  
    & & & \multicolumn{2}{p{0.3\linewidth}}{
        \cellcolor[rgb]{0.9,0.9,0.9}{
            \makecell[{{p{\linewidth}}}]{
                \texttt{\tiny{[GM$|$GM]}}
                \texttt{Zdi se mi, da imamo drugačne priority, kar je odlično za dogovor. Mene najbolj zanima **prstan** (vrednost 5), nato pa **ura** (vrednost 3). **Klešče** imajo za mene vrednost 1, **opeke** pa 0.} \\
\\ 
\texttt{Predlagam sledeč razdelek:} \\
\texttt{*   **Jaz** vzamem **prstan** in **eno opeko** (opeka mi je vseeno, a jo vzamem, da uravnotežim število predmetov).} \\
\texttt{*   **Ti** vzameš **uro** in **klešče** (torej tvoja dva najvrednejša predmeta).} \\
\\ 
\texttt{Tako bi vsak od nas dobil to, kar mu je najbolj pomembno. Kaj misliš?} \\
            }
        }
    }
    & & \\ \\

    \theutterance \stepcounter{utterance}  
    & & \multicolumn{4}{p{0.6\linewidth}}{
        \cellcolor[rgb]{0.9,0.9,0.9}{
            \makecell[{{p{\linewidth}}}]{
                \texttt{\tiny{[P1$\langle$GM]}}
                \texttt{Zdi se mi, da imamo drugačne priority, kar je odlično za dogovor. Mene najbolj zanima **prstan** (vrednost 5), nato pa **ura** (vrednost 3). **Klešče** imajo za mene vrednost 1, **opeke** pa 0.} \\
\\ 
\texttt{Predlagam sledeč razdelek:} \\
\texttt{*   **Jaz** vzamem **prstan** in **eno opeko** (opeka mi je vseeno, a jo vzamem, da uravnotežim število predmetov).} \\
\texttt{*   **Ti** vzameš **uro** in **klešče** (torej tvoja dva najvrednejša predmeta).} \\
\\ 
\texttt{Tako bi vsak od nas dobil to, kar mu je najbolj pomembno. Kaj misliš?} \\
            }
        }
    }
    & \\ \\

    \theutterance \stepcounter{utterance}  
    & \multicolumn{4}{p{0.6\linewidth}}{
        \cellcolor[rgb]{0.9,0.9,0.9}{
            \makecell[{{p{\linewidth}}}]{
                \texttt{\tiny{[P1$\rangle$GM]}}
                \texttt{Zvendi se mi odličen predlog! Tako bi jaz dobil uro (7) in klešče (1) = **8 točk**, ti pa prstan (5) in opeko (0) = **5 točk**. Skupaj dosežemo **13 točk**, kar je najboljši možni skupni rezultat (maksimalna vsota vrednosti vseh predmetov).} \\
\\ 
\texttt{Opeka, ki jo vzameš ti, mi ni potrebna (vrednost 0), tako da je to idealna rešitev za oba.} \\
\\ 
\texttt{Strinjam se s predlogom.} \\
\\ 
\texttt{{[}Predlog: 1 ura, 1 klešče{]}} \\
            }
        }
    }
    & & \\ \\

    \theutterance \stepcounter{utterance}  
    & & & \multicolumn{2}{p{0.3\linewidth}}{
        \cellcolor[rgb]{0.9,0.9,0.9}{
            \makecell[{{p{\linewidth}}}]{
                \texttt{\tiny{[GM$|$GM]}}
                \texttt{proposal} \\
            }
        }
    }
    & & \\ \\

    \theutterance \stepcounter{utterance}  
    & & & \multicolumn{2}{p{0.3\linewidth}}{
        \cellcolor[rgb]{0.9,0.9,0.9}{
            \makecell[{{p{\linewidth}}}]{
                \texttt{\tiny{[GM$|$GM]}}
                [1, 0, 0, 1]
            }
        }
    }
    & & \\ \\

    \theutterance \stepcounter{utterance}  
    & & \multicolumn{4}{p{0.6\linewidth}}{
        \cellcolor[rgb]{0.9,0.9,0.9}{
            \makecell[{{p{\linewidth}}}]{
                \texttt{\tiny{[GM$\rangle$P2]}}
                \texttt{Drugi igralec je predložil skrivni predlog. Zdaj ste na vrsti, da predložite svoj predlog.} \\
            }
        }
    }
    & \\ \\

    \theutterance \stepcounter{utterance}  
    & & & \multicolumn{4}{p{0.6\linewidth}}{
        \cellcolor[rgb]{0.9,0.9,0.9}{
            \makecell[{{p{\linewidth}}}]{
                \texttt{\tiny{[GM$\langle$P2]}}
                \texttt{{[}Predlog: 1 ura, 1 klešče{]}} \\
            }
        }
    }
     \\ \\

    \theutterance \stepcounter{utterance}  
    & & & \multicolumn{2}{p{0.3\linewidth}}{
        \cellcolor[rgb]{0.9,0.9,0.9}{
            \makecell[{{p{\linewidth}}}]{
                \texttt{\tiny{[GM$|$GM]}}
                \texttt{proposal} \\
            }
        }
    }
    & & \\ \\

    \theutterance \stepcounter{utterance}  
    & & & \multicolumn{2}{p{0.3\linewidth}}{
        \cellcolor[rgb]{0.9,0.9,0.9}{
            \makecell[{{p{\linewidth}}}]{
                \texttt{\tiny{[GM$|$GM]}}
                [1, 0, 0, 1]
            }
        }
    }
    & & \\ \\

    \theutterance \stepcounter{utterance}  
    & & & \multicolumn{2}{p{0.3\linewidth}}{
        \cellcolor[rgb]{0.9,0.9,0.9}{
            \makecell[{{p{\linewidth}}}]{
                \texttt{\tiny{[GM$|$GM]}}
                [[1, 0, 0, 1], [1, 0, 0, 1]]
            }
        }
    }
    & & \\ \\

\end{supertabular}
}

\end{document}
